\section{FORMAL RESTATEMENT}
\textbf{Erd\H{o}s \#133; reconstruction, 2026-09-23.}
Let $f(n)$ denote the smallest possible maximum degree of an $n$-vertex
triangle-free graph of diameter $2$. Determine its order of growth;
must $f(n)/\sqrt n\to\infty$? We establish $f(n)=\Theta(\sqrt n)$
using an explicitly stated theorem on complete sum-free sets.

\section{QUICK LITERATURE AND CONTEXT CHECK}
The literal word ``minimal'' in the webpage's universal-threshold
phrasing would give a trivial value; the extremal interpretation above
is the one used in its bounds and references. The negative answer and
order of growth are established results. We prove the degree lower
bound, the Cayley-graph translation, and a further explicit construction.
For the sharp order upper bound for every large $n$, we use Haviv--Levy's
Theorem 1.5 as an external input. Its combinatorial construction is not
claimed as proved here.

\section{WORK}
\subsection*{The unavoidable degree scale}
If a graph has maximum degree $d$ and diameter $2$, choose a vertex $v$.
There are at most $d$ vertices at distance one, and each has at most
$d-1$ other neighbors. Therefore
\[
 n\le1+d+d(d-1)=d^2+1,\qquad f(n)\ge\sqrt{n-1}.           \tag{1}
\]
No assumption that the second neighborhoods are disjoint is needed:
overcounting is in the correct direction.

\subsection*{An explicit intermediate construction}
For $m\ge2$, take the $m$-subsets of $[3m-1]$ as vertices and join
disjoint sets. Three pairwise disjoint $m$-sets would require $3m$
ground elements, so the graph is triangle-free. Any two different
intersecting vertices have union of size at most $2m-1$, leaving at
least $m$ elements for a common disjoint neighbor. Disjoint vertices
are already adjacent. The diameter is therefore $2$, and the graph is
regular with
\[
 n_m=\binom{3m-1}{m},\qquad d_m=\binom{2m-1}{m}.
\]
Stirling's formula gives
$\log n_m=(3\log3-2\log2)m+O(\log m)$ and
$\log d_m=2m\log2+O(\log m)$, whence
\[
 d_m=n_m^{\,2\log2/(3\log3-2\log2)+o(1)}.
\]
This verifies the historical exponent approximately $0.726$, rather
than assuming that this construction already has square-root degree.

\subsection*{The full order bound from a sum-free input}
We now use Haviv--Levy's theorem: there is an absolute $C$ such that,
for every sufficiently large $n$, a set $S\subseteq\mathbb Z/n\mathbb Z$
exists with
\[
 S=-S,\quad S\cap(S+S)=\varnothing,\quad
 S\cup(S+S)=\mathbb Z/n\mathbb Z,\quad |S|\le C\sqrt n.   \tag{2}
\]
The word ``complete'' in this theorem means exactly the third
condition, not that $S$ is a complete interval or a complete graph.
Define a graph on the group by joining $x$ and $y$ when $y-x\in S$.
Condition (2) forces $0\notin S$: otherwise $0\in S+S$. Symmetry makes
the graph undirected, and every vertex has exactly $|S|$ neighbors.
A triangle $x,y,z$ would give $y-x,z-y\in S$ and
$z-x=(y-x)+(z-y)\in S$, contradicting sum-freeness.
If $x\ne y$ are not adjacent, their difference belongs to $S+S$.
Write $y-x=s+t$ with $s,t\in S$; then $x,x+s,y$ is a two-edge path,
with distinct intermediate vertex because $s,t\ne0$. Thus the graph
has diameter at most two, and for large $n$ it is not complete, so
its diameter is two. Hence
\[
 f(n)\le|S|\le C\sqrt n.                                \tag{3}
\]
Together (1) and (3) prove the claimed order and disprove divergence.

\section{VERIFICATION}
The intermediate subset graph is checked directly for small $m$.
As a small Cayley example, $S=\{1,4\}\subseteq\mathbb Z/5\mathbb Z$
satisfies (2) and gives the cycle $C_5$, of degree two and diameter two.
These examples illustrate the translation; they do not replace the
all-$n$ existence theorem. The numerical exponent follows the displayed
exact logarithmic expression, independently of decimal transcription
in a secondary source.

\section{FINAL STATUS}
\textbf{SOLVED (known order of growth deduced from an explicit external theorem).}
The graph-theoretic deduction is complete. Haviv--Levy Theorem 1.5 is
an essential omitted proof dependency. Precise asymptotics for
$f(n)/\sqrt n$ are not asserted; no new result is claimed.

\section{SOURCES}
\url{https://www.erdosproblems.com/133}.
I.~Haviv and D.~Levy, \emph{Symmetric complete sum-free sets in cyclic
groups}, Theorem 1.5 and Section 1.2,
\url{https://arxiv.org/abs/1703.04118}.
Checked 2026-09-23. Completion is relative to (2), not a claim to
reproduce its proof or settle the optimal leading constant.

\section{COMPLETION ESTIMATE}
\textbf{COMPLETION: 100\%}
